\documentclass[10pt,a4paper]{amsart}
\usepackage[margin=19mm]{geometry}
\usepackage{amsmath,amssymb,mathtools}
\usepackage[scr=rsfs,cal=euler]{mathalfa}
\usepackage{xfrac}
\usepackage[unicode,hidelinks,bookmarks=false]{hyperref}
\newcommand{\ie}{i.e.\ }
\newcommand{\cf}{cf.\ }
\newcommand{\resp}{respectively\ }
\newcommand{\Subsection}{\subsection}
\providecommand{\nobreakdash}{\nobreak\hbox{-}}
\setlength{\emergencystretch}{2em}
\multlinegap0pt
\usepackage{siunitx}
\usepackage{siunitx}
\usepackage{siunitx}
\usepackage{siunitx}
\usepackage{xcolor}
\usepackage[shortlabels]{enumitem}
\usepackage{amssymb}
\usepackage{algorithm}\usepackage{algpseudocode}
\usepackage{siunitx}
\usepackage{bm}
\newtheorem{assumption}{Assumption}
\newtheorem{theorem}{Theorem}
\newtheorem{proposition}{Proposition}
\def\IE{\mathbb{E}}
\def\R{\mathbb{R}}
\def\bC{\mathbf{C}}
\def\bT{\mathbf{\Theta}}
\newcommand{\minparam}{\theta^\star_K}
\title{Sharp Gaussian approximations for Decentralized Federated Learning}
\author{Soham Bonnerjee, Sayar Karmakar, Wei Biao Wu}
\date{Source: arXiv:2505.08125v4 (CC BY 4.0)}
\begin{document}
\maketitle

\noindent\emph{Source and licence.} Sharp Gaussian approximations for Decentralized Federated Learning. Version arXiv:2505.08125v4 (CC BY 4.0), distributed by arXiv under the Creative Commons Attribution 4.0 International licence, http://creativecommons.org/licenses/by/4.0/, as recorded in the arXiv licence metadata for this exact version and shown on the abstract page https://arxiv.org/abs/2505.08125v4. Full licence notice: \texttt{licenses/arxiv-2505.08125-CC-BY-4.0.txt}.\par\smallskip

\noindent\textit{Editorial scope.} Counted unit: the article's proposition prop:fourth-moment-original (source0.tex lines 1027--1032). The article's complete original proof of it is delivered with it (source0.tex lines 1189--1236, one \texttt{proof} environment of 48 lines, which the article places away from the statement). No mathematics is written, repaired or re-derived: the statement and the proof are the article's own lines, in the article's own notation, and no retained line is substituted or re-flowed.  Retained uncounted as the source-local context the proof needs: lines 384--386; lines 393--395; lines 397--405; lines 611--614; lines 616--619; lines 621--623; lines 1020--1025.  Nothing was omitted from any retained span, so the package carries no omission record.  No retained line contains a citation, so the wrapper carries no bibliography and no bibliography entry.  No retained line contains a figure environment, an image-inclusion command or drawing code, so no image is delivered and the package declares no asset dependency.  Editorial formatting only, in addition: an AMS \texttt{amsart} preamble that restates the article's own declarations and macros used by the retained lines, namely the environments \texttt{assumption}, \texttt{theorem}, \texttt{proposition} on separate counters, together with the standard packages those lines need.  The retained environments print in this document as Assumption 1, Theorem 1, Proposition 1 in order of appearance, whereas the article prints the same items with the numbering of its own full document.  The complete native source of the article as distributed in the arXiv e-print for this exact version is bundled unchanged as \texttt{sources/178016/source0.tex}.
\begin{assumption}\label{ass:sync}
    The connection matrix $\bC$ satisfies $\bC \mathbf{1}=\mathbf{1}$ and $\bC^\top=\bC$. Moreover, if $\lambda_1\geq \ldots \geq \lambda_K$ denote the ordered eigen-values of $\bC$, then $\lambda_1=1$, and $\lambda_2 = \rho <1$ for some $\rho \in (0,1)$.
\end{assumption}

\begin{align} \label{eq:Y_bar_defn}
    \Bar{Y}_n:= n^{-1}\sum_{t=1}^n Y_t = K^{-1} \sum_{k=1}^K n^{-1}\sum_{t=1}^n \theta_t^k,
\end{align}

\begin{theorem}\label{thm: berry-esseen}
   Define $\mathcal{A}_s^t:=\prod_{j=s+1}^t(I-\eta_tA)$, $\mathcal{A}_t^t=I$, where $A:=\nabla_2 F(\minparam)$ for $t \in [n]$. Further, for $s\in [n]$, define the random vectors 
   \[ u_s = \eta_s\sum_{k=1}^K w_k\bigg(\sum_{j=s}^n \mathcal{A}_s^j\bigg) g_k(\minparam, \xi_s^k), \text{with $\Sigma_n:= n^{-1}\sum_{s=1}^n\IE[|u_s u_s^\top|]$, $g_k(\theta, \xi^{k}) = \nabla F_k(\theta) - \nabla f_k(\theta, \xi^k)$.} \]
Let there exist a constant $C$ such that for $\xi^k\sim \mathcal{P}_k, k\in [K]$, it holds  $\max_{k \in [K]} \IE[|g_k(\minparam, \xi^k)|^2]\leq C$. Suppose that the step-size schedules of the clients satisfy that $ \eta_{t} = \eta_0(t+ k_0)^{-\beta}$ for some fixed $\eta_0, k_0>0,$ and $\beta \in (1/2, 1)$. Then, under Assumptions \ref{ass:sync}, \ref{ass:sc} and \ref{ass:Lipschitz} and \ref{ass:boundedloss} with $p=4$, and $\bar{Y}_n$ as in \eqref{eq:Y_bar_defn}, it holds that
    \begin{align}\label{eq:berry-esseen}
        d_{\mathrm{C}}( \sqrt{n}(\bar{Y}_n - \minparam) , Z)  \lesssim \frac{1}{\sqrt{nK}}  + n^{\frac{1}{2}-\beta} \sqrt{K} + \frac{n^{-\frac{\beta}{2}}}{\sqrt{K}} ,
    \end{align}
    where $\lesssim$ hides constants involving $d, \beta, \mu, L$ and $\rho$, and $Z \sim N(0, \Sigma_n)$. 
\end{theorem}

\begin{assumption}[Strong convexity]\label{ass:sc}
    There exists $\mu>0$ such that for each $k\in [K]$, 
    \[ \langle\nabla F_k(\theta) - \nabla F_k(\theta'), \theta-\theta'\rangle \geq \mu |\theta-\theta'|^2, \ \theta, \theta' \in \R^d. \]
\end{assumption}

\begin{assumption}(Stochastic Lipschitzness of noisy gradients)\label{ass:Lipschitz}
    There exists $L>0$ such that for each $k \in [K]$,
    \[ \IE_{\mathcal{P}_k}[|\nabla f_k(\theta, \xi^k)- f_k(\theta', \xi^k)|^2] \leq L|\theta - \theta'|^2, \ \theta, \theta' \in \R^d.\]
\end{assumption}

\begin{assumption}(Control on noisy gradients)\label{ass:boundedloss}
The functions $f_k(\theta, \xi)$ is assumed to be continuously differentiable with respect to $\theta$ for any fixed $\xi$. Moreover, assume that $\max_{k\in [K]}\IE[|g_k(\theta, \xi^k)|^p]< \infty$ for some $p \geq 2$.
\end{assumption}

\begin{proposition}\label{prop:fourth-moment-difference}
  Under the conditions of Theorem \ref{thm: berry-esseen}, it holds that 
    \begin{align} \label{eq:fourth-moment-difference}
        \IE[|\bT_t(I-J)|_F^4]=O(\eta_t^4K^2).
    \end{align}
\end{proposition}

\begin{proposition}\label{prop:fourth-moment-original}
   Under the conditions of Theorem \ref{thm: berry-esseen}, it holds that
    \begin{align} \label{eq:fourth-moment-original}
        \IE[|Y_t - \minparam|^4]=O(\frac{\eta_t^2}{K^2} + \eta_t^4).
    \end{align}
\end{proposition}

\begin{proof}[Proof of Proposition \ref{prop:fourth-moment-original}]
    Write  
\begin{align}
    R_t:=Y_t - \minparam &= E_1 + E_2 + E_3, \text{ where},\nonumber\\
    E_1 & = R_{t-1} - \eta_t\nabla F(Y_{t-1}) , \nonumber\\
    E_2 &= \eta_t \sum_{k=1}^K w_k (\nabla F_k(Y_{t-1}) - \nabla F_k(\theta_{t-1}^k)),  \text{ and} \nonumber\\
    E_3 &= \eta_t \sum_{k=1}^K w_k g_k(\theta_{t-1}^k, \xi_t^k). \label{eq:foruth-moment-decomp}
\end{align}
Note that trivially, Assumptions \ref{ass:sc} and \ref{ass:Lipschitz} imply that 
\begin{align}
    \IE[|E_1|^4] \leq (1- \eta_t c)\IE[|R_{t-1}|^4] \label{eq:E1}.
\end{align}
Moving on, for $E_2$ we proceed as follows:
\begin{align}
    \IE[|E_2|^4] = & \eta_t^4 \IE[|\sum_{k=1}^K w_k(\nabla F_k(Y_{t-1}) - \nabla F_k(\theta_{t-1}^k))|^4] \nonumber\\
    \leq & C_2\frac{\eta_t^4}{K^2} \IE\bigg[ \Big( \sum_{k=1}^K |Y_{t-1} - \theta_{t-1}^k|^2 \Big)^2 \bigg] \nonumber\\
    \leq & C_2\frac{\eta_t^4}{K^2} \IE[|\Theta_{t-1}(I-J)|_F^4] \leq C_2'{\eta_t^8} \label{eq:E2},
\end{align}
for some constants $C_2, C_2'>0$, where the final assertion is drawn from Proposition \ref{prop:fourth-moment-difference}. Finally, for $E_3$, we obtain,
\begin{align}
    \IE[|E_3|^4] = & \eta_t^4 \IE[|\sum_{k=1}^K w_k g_k(\theta_{t-1}^k, \xi_t^k)|^4] \leq  C_3 \frac{\eta_t^4}{K^2} \label{eq:E3},
\end{align}
 for some constant $C_3>0$, where we have used the fact that $g_k(\theta_{t-1}^k, \xi_t^k)$ are mean-zero and independent random vectors conditional on $\mathcal{F}_t$. Now, we will leverage \eqref{eq:E1}-\eqref{eq:E3} to develop bounds on the cross-product terms. In particular,
\begin{align}
    \IE[|E_1|^2 |E_2|^2] \leq & \sqrt{\IE[|E_1|^4]} \sqrt{\IE[|E_l|^4]} \nonumber\\
    \leq  & C_3 \eta_t^4 \sqrt{\IE[|E_1|^4]} \nonumber\\
    \leq & C_3 \min\{\eta_t \varepsilon \IE[|E_1|^4] + (4\varepsilon)^{-1} \eta_t^7 , \  \eta_t^2 \varepsilon \IE[|E_1|^4] + (4\varepsilon)^{-1}  \eta_t^6\}, \label{eq:E1-cross-prods-1}
\end{align}
and similarly
\begin{align}
    \IE[|E_1|^2 |E_3|^2] \leq & C_3 \min\{\eta_t \varepsilon \IE[|E_1|^4] + (4\varepsilon)^{-1} \frac{\eta_t^3}{K^2} , \  \eta_t^2 \varepsilon \IE[|E_1|^4] + (4\varepsilon)^{-1} \frac{\eta_t^2}{K^2} \}, \label{eq:E1-cross-prods-2}
\end{align}
where the final assertions follow from Young's inequality. Here, $\varepsilon$ is chosen to be small enough, however it remains a constant; the explicit choice of $\varepsilon$ will be indicated towards the end of the proof, when we collect terms to establish the recursion. Note that, quite trivially, from \eqref{eq:E2} and \eqref{eq:E3}, one has
\begin{align}
    \IE[|E_2|^2 |E_3|^2] \leq C_4\frac{\eta_t^6}{K} \text{ for some constant $C_4>0$}. \label{eq:E2-E3}
\end{align}
Rest of the cross-products are strictly dominated by some combinations of the terms analyzed till now. For example, for $l,r, q \in \{1,2,3\}$, Cauchy-Schwarz and AM-GM inequalities implies that
\begin{align}
  \IE[(E_l^\top E_q)^2] \leq & \IE[|E_l|^2 |E_q|^2], \nonumber\\
    \IE[|E_l|^2 (E_r^\top E_q) ] \leq & \sqrt{\IE[|E_l|^4]} \sqrt{\IE[(E_r^\top E_q)^2]}, \nonumber\\
    \IE[|(E_l^\top E_r) (E_l^\top E_q)|] \leq & 2^{-1}(\IE[(E_l^\top E_r)^2] + \IE[(E_l^\top E_q)^2]) \label{eq:rest-of-cross-prod}.
\end{align}
A careful collection of terms from \eqref{eq:E1}-\eqref{eq:rest-of-cross-prod} yields
\begin{align}
    \IE[|R|_t^4]\leq (1-\eta_tc) (1+ \eta_t C_0 \varepsilon)\IE[|R_{t-1}|^4] + O(\frac{\eta_t^3}{K^2} + \eta_t^5),\label{eq:final R_t}
\end{align}
where $C_0$ is a large constant depending upon $C_2, C_3$ and $C_4$. Now, choose $\varepsilon>0$ so that $C_0\varepsilon < c/2$, upon which we immediately obtain $(1-\eta_tc) (1+ \eta_t C_0 \varepsilon)< 1- \eta_t c/2$. Therefore, \eqref{eq:final R_t} immediately yields \eqref{eq:fourth-moment-original}. 
\end{proof}

\end{document}
